\begin{center}
    \textbf{\Huge{Relatório - Síntese e Decisão Técnica: Estratégias de Execução de Modelos de IA}}
\end{center}

\section{Introdução}

O propósito deste documento é consolidar os resultados das pesquisas conduzidas nos estudos
técnicos associados aos Cards S0-3A a S0-3D do projeto MIIA, que investigaram diferentes abordagens
para execução e integração de modelos de inteligência artificial em ambientes de simulação
construtiva.

Os estudos analisaram quatro estratégias principais:

\begin{itemize}
\item Execução por meio de runtimes nativas (LibTorch e TensorFlow C++ API);
\item Uso do padrão Open Neural Network Exchange (ONNX) como formato intermediário de modelos;
\item Servimento remoto de modelos (model serving);
\item Embedding do interpretador Python em aplicações C++.
\end{itemize}

O objetivo desta síntese é comparar essas abordagens sob a perspectiva de desempenho,
complexidade de integração, riscos técnicos, esforço de manutenção e aderência ao contexto do
projeto, culminando na seleção de soluções candidatas para implementação em provas de conceito
(PoC) ou MVP nas próximas etapas do projeto.

\subsection{Referências}

\begin{itemize}
    \item Documento n\textsuperscript{o} MIIA-REL-46ff681a, \textbf{Relatório - Estudo de Runtimes Nativas (Libtorch e TensorFlow C++ API)}, Instituto de Estudos Avançados (IEAv), 2025.
    \item Documento n\textsuperscript{o} MIIA-REL-323690a5, \textbf{Relatório - Model Serving: Práticas e Solucões existentes}, Instituto de Estudos Avançados (IEAv), 2025.
    \item Documento n\textsuperscript{o} MIIA-REL-333869a8, \textbf{Relatório - Open Neural Network Exchange (ONNX) e Possíveis Aplicações}, Instituto de Estudos Avançados (IEAv), 2025.
    \item Documento n\textsuperscript{o} MIIA-REL-d6ee0363, \textbf{Relatório - Pesquisa de Embedding de Python em C++}, Instituto de Estudos Avançados (IEAv), 2025.
\end{itemize}

\section{Síntese dos Estudos}

Esta seção resume os principais achados dos estudos técnicos conduzidos nos CARDS S0-3A a S0-3D,
os quais investigaram diferentes estratégias para integração e execução de modelos de inteligência
artificial em sistemas de simulação construtiva.

Os resultados indicam que as abordagens avaliadas atendem a diferentes paradigmas de decisão
computacional, que podem coexistir dentro da arquitetura do projeto.

\subsection{Runtimes nativas (LibTorch e TensorFlow C++ API)}

Runtimes nativas permitem executar modelos diretamente em aplicações C++, eliminando a
dependência de ambientes interpretados ou serviços externos.

A inferência ocorre dentro do próprio processo da aplicação, reduzindo a latência e permitindo
integração direta com sistemas legados. Estudos indicam latências na ordem de milissegundos,
adequadas para aplicações que exigem execução em tempo quase real.

Entretanto, essa abordagem introduz desafios relevantes de engenharia, incluindo aumento
significativo do tamanho dos binários, maior complexidade no processo de build e manutenção
de dependências pesadas de frameworks de aprendizado de máquina.

Além disso, mudanças frequentes nas APIs e incompatibilidades podem aumentar o custo
de manutenção ao longo do tempo.

No contexto do projeto, observou-se que o uso direto dessas runtimes tende a se tornar
desnecessário quando os modelos podem ser exportados para um formato intermediário
padronizado, como o ONNX.


\subsection{ONNX como formato intermediário}

O ONNX foi analisado como um mecanismo de interoperabilidade entre frameworks de treinamento
e motores de inferência.

Seu principal benefício consiste na separação entre treinamento, representação do modelo e
execução, permitindo que modelos treinados em diferentes frameworks sejam implantados em
ambientes heterogêneos.

Essa separação reduz o acoplamento arquitetural e facilita a evolução independente dos
componentes do sistema, sendo particularmente relevante em sistemas de longa vida útil,
como simuladores complexos.

Além disso, modelos baseados em redes neurais — incluindo modelos utilizados em
Reinforcement Learning — podem ser exportados para ONNX e executados por runtimes
especializados sem dependência direta do framework de treinamento.

Dessa forma, o ONNX pode atuar como um contrato arquitetural entre o ambiente de treinamento
de modelos e o simulador, reduzindo a necessidade de integração direta com frameworks como
PyTorch ou TensorFlow.


\subsection{Servimento remoto de modelos}

O modelo de servimento consiste em disponibilizar modelos de IA como serviços acessados via
rede por meio de APIs ou microserviços.

Essa abordagem facilita escalabilidade, versionamento e gerenciamento centralizado de modelos,
características comuns em arquiteturas modernas de ML em produção.

Entretanto, em ambientes de simulação construtiva com requisitos de latência previsível e
execução síncrona, a dependência de comunicação em rede pode introduzir variabilidade
temporal e comprometer a consistência do ciclo de simulação.

Assim, embora seja adequada para aplicações distribuídas, essa abordagem apresenta limitações
para integração direta em loops de simulação de alta frequência.


\subsection{Embedding de Python em C++}

O embedding do interpretador Python permite executar código Python diretamente dentro de uma
aplicação C++, possibilitando acesso ao ecossistema de bibliotecas disponíveis nessa linguagem.

Essa abordagem permite que scripts ou módulos Python sejam utilizados para implementar
comportamentos complexos ou experimentais, funcionando como um mecanismo de extensibilidade
do sistema.

Entretanto, essa integração introduz custos operacionais adicionais, como gerenciamento do
interpretador, conversão de dados entre linguagens e controle do Global Interpreter Lock (GIL).

Apesar dessas limitações, o embedding Python apresenta valor arquitetural relevante ao
possibilitar a implementação de lógicas procedurais flexíveis, complementando outras
abordagens de tomada de decisão utilizadas no simulador.

\section{Matriz Comparativa das Alternativas}

\begin{center}
\renewcommand{\arraystretch}{1.4}

\begin{tabular}{|p{3cm}|p{2.5cm}|p{3cm}|p{3.5cm}|p{3cm}|}
\hline
\textbf{Abordagem} &
\textbf{Desempenho Esperado} &
\textbf{Complexidade de Integração} &
\textbf{Principais Riscos} &
\textbf{Aderência ao Projeto} \\
\hline

Runtimes nativas
(LibTorch / TensorFlow C++)

& Alta

& Alta

& Dependências pesadas de framework,
complexidade de build e risco de incompatibilidades de API

& Baixa a moderada.
Tornam-se desnecessárias quando modelos podem ser exportados para ONNX. \\

\hline

ONNX + Runtime de Inferência

& Alta

& Média

& Limitações ocasionais de operadores ou conversão de modelos

& Alta.
Permite desacoplamento entre treinamento e execução e atende
modelos baseados em redes neurais e RL. \\

\hline

Model Serving

& Média

& Média

& Latência de rede e variabilidade temporal

& Baixa para integração direta em loops de simulação síncronos. \\

\hline

Embedding Python em C++

& Média

& Média

& Overhead do interpretador,
contenção do GIL e complexidade de integração

& Alta para lógica procedural e extensibilidade do sistema. \\

\hline

\end{tabular}
\end{center}

\section{Decisão Técnica}

Com base nos resultados consolidados dos estudos, definiu-se uma arquitetura que contempla
três paradigmas complementares de tomada de decisão no simulador.

\begin{itemize}

\item Modelos baseados em aprendizado de máquina (incluindo Reinforcement Learning),
executados por meio de modelos exportados no formato ONNX;

\item Estruturas de decisão baseadas em árvores de comportamento (Behavior Trees),
mantidas por meio da biblioteca BehaviorTree.cpp já utilizada no projeto;

\item Lógicas procedurais e comportamentos experimentais implementados por meio
de scripts ou módulos Python integrados via embedding do interpretador Python em C++.

\end{itemize}

Nesse modelo arquitetural, o ONNX atua como contrato para integração de modelos de
aprendizado de máquina, permitindo desacoplamento entre o ambiente de treinamento
e o ambiente de execução do simulador.

As Behavior Trees permanecem como mecanismo estruturado de composição de decisões,
utilizando a infraestrutura já existente baseada em XML.

Por fim, o embedding Python fornece um mecanismo de extensibilidade capaz de suportar
lógicas procedurais ou experimentais que não se encaixem diretamente nas duas
abordagens anteriores.

O uso direto de runtimes nativas como LibTorch ou TensorFlow C++ API é considerado
desnecessário para a arquitetura principal do sistema, sendo reservado apenas para
cenários excepcionais em que a exportação para ONNX não seja viável.

Com esta decisão, considera-se concluída a fase exploratória dos CARDS S0-3A a S0-3D,
habilitando o início das implementações experimentais nas próximas sprints.
